% ---------------------------------------------------------------------------
% Day 6 review: this frame used to show the SAC v2 loop with BOTH the
% automatic-temperature line and the clipped double-Q highlighted, before
% either had been motivated. It now shows the base algorithm with a FIXED
% alpha; the single extra v2 line is introduced in auto_temp_derivation.tex,
% straight after the dual derivation that justifies it.
% ---------------------------------------------------------------------------
\begin{frame}{SAC: Algorithm}
\begin{center}
{\scriptsize\textbf{\textcolor{red}{Red}} $=$ the \textbf{two} things that differ from TD3 (Day 5): the actor is \emph{stochastic}, so actions are \textbf{sampled} from it, and the TD target carries the \emph{entropy term}.\\
Shorthand: $Q^{\min}_{\phi^{-}}$ means $\min_{i=1,2} Q_{\phi^{-}_i}$.\quad $\alpha$ is a \textbf{fixed hyperparameter} here.\quad
$\lambda_Q,\lambda_\pi$ are \textbf{learning rates} (Adam step sizes) for critic and actor; $\tau$ is the target-network mixing rate from Day 5.}
\end{center}
\begin{center}
\resizebox{\textwidth}{!}{%
\begin{minipage}{1.06\textwidth}
\begin{algorithm}[H]
\scriptsize
\DontPrintSemicolon
\LinesNotNumbered
\setlength{\algomargin}{0pt}
Initialize actor $\theta$, critics $\phi_1,\phi_2$;\ \ targets $\phi^{-}_i\!\leftarrow\!\phi_i$;\ \ replay $\mathcal{D}\leftarrow\emptyset$;\ \ choose $\alpha$\;
\For{each iteration}{
  \For{each environment step}{
    \textcolor{red}{$\mathbf{a}_t\sim\pi_\theta(\,\cdot\mid\mathbf{s}_t)$}\ \ \emph{(sampled --- no injected noise)};\ \ $\mathbf{s}_{t+1}\sim p$;\ \
    $\mathcal{D}\leftarrow\mathcal{D}\cup\{(\mathbf{s}_t,\mathbf{a}_t,r_t,\mathbf{s}_{t+1})\}$\;
  }
  \For{each gradient step}{
    sample minibatch from $\mathcal{D}$;\ \ \textcolor{red}{$\tilde{\mathbf{a}}_{t+1}\sim\pi_\theta(\,\cdot\mid\mathbf{s}_{t+1})$}\ \ \emph{(sampled from the current actor)}\;
    $y = r_t + \gamma\big( Q^{\min}_{\phi^{-}}(\mathbf{s}_{t+1},\tilde{\mathbf{a}}_{t+1}) \textcolor{red}{\;-\; \alpha\log\pi_\theta(\tilde{\mathbf{a}}_{t+1}\mid\mathbf{s}_{t+1})}\big)$\;
    $\phi_i \leftarrow \phi_i - \lambda_Q\,\hat\nabla_{\phi_i}\big[\tfrac12(Q_{\phi_i}(\mathbf{s}_t,\mathbf{a}_t)-y)^2\big]$,\ \ $i\in\{1,2\}$\;
    $\theta \leftarrow \theta - \lambda_\pi\,\hat\nabla_\theta J_\pi(\theta)$ \ \ ($J_\pi$ critic term uses $Q^{\min}_{\phi}$)\;
    $\phi^{-}_i \leftarrow \tau\phi_i+(1-\tau)\phi^{-}_i$,\ \ $i\in\{1,2\}$\;
  }
}
\end{algorithm}
\end{minipage}%
}
\end{center}
\vspace{-0.2em}
{\scriptsize Black $=$ Day 5 unchanged: replay buffer, twin critics, clipped double-$Q$ target, Polyak target critics. \textbf{But the actor is not TD3's} --- it outputs a distribution $(\mu_\theta,\sigma_\theta)$, never a bare action, which is why SAC needs \emph{no target actor} and \emph{no smoothing noise}: TD3 needed both only because its actor was deterministic.}
\end{frame}
